% =====================================================================
%  Climbing Wall Force-Measurement System
%  Calibration Guide
%
%  Compile with:   latexmk -pdf -synctex=1 calibration_guide.tex
%  Screenshots live in ./images/calibration_guide/
% =====================================================================
\documentclass[11pt,a4paper]{article}
\usepackage{climbingwall}
\graphicspath{{images/calibration_guide/}}
\setdocname{Calibration Guide}
\renewcommand{\shotplacement}{H}  % pin the walkthrough screenshots to their steps

\begin{document}

\guidetitle{Calibrating the Force-Measurement System}%
  {Step-by-step guide --- how and why the sensors are calibrated.}

% =====================================================================
\section{Why do we calibrate?}
% =====================================================================

\subsection{What the sensors actually measure}
Each of the four boards contains a \textbf{three-axis load cell}: a small device
that produces a tiny electrical signal when it is pushed or pulled. The boards are
plugged into the computer over USB, and the backend reads all four about 100 times
per second, delivering twelve numbers per reading to the dashboard (four boards
$\times$ three directions, $X$, $Y$ and $Z$).
The catch is that these raw numbers are \textbf{voltage ratios}, not forces. They
are extremely small (around $10^{-6}$) and every individual sensor is slightly
different. A reading of, say, $0.000012$ from one board does not directly tell you
how many newtons are pressing on it. We need a rule that converts the raw number
into a force. \textbf{Finding that rule, for each sensor, is what calibration
does.}

\subsection{The conversion rule}
For every sensor channel we assume a simple \emph{straight-line} relationship
between the raw reading and the true force:
\begin{equation}
  \text{force (N)} \;=\; \bigl(\text{raw reading} - \text{offset}\bigr)\times\text{scale}.
  \label{eq:line}
\end{equation}
Two unknowns appear here, and they behave \emph{very differently} --- which is why
the system handles them in two separate ways:
\begin{itemize}[leftmargin=1.4em]
  \item \textbf{scale} --- how many newtons correspond to one unit of raw signal
        (the steepness, or slope, of the line). This is a \emph{fixed physical
        property} of the sensor; it barely changes over time. You measure it
        \textbf{once}, with known weights, and it is \textbf{saved}. Finding the
        scale is what \textbf{calibration} means in this guide.
  \item \textbf{offset} --- the raw number the sensor reports when \emph{nothing} is
        loaded (its ``zero point''). This one \emph{drifts} --- with temperature, as
        the rig settles, whenever the hardware is re-seated --- so a value measured
        today is wrong tomorrow. Rather than save it, you re-set it \textbf{live, on
        demand}, with the \textsf{Zero Sensors} button (a ``tare''). The offset is
        \textbf{never saved} as part of the calibration.
\end{itemize}
In short: the \textbf{scale} is stable, so it is calibrated and stored; the
\textbf{zero} drifts, so it is tared fresh whenever you need it
(Section~\ref{sec:zeroing}).

\subsection{Three things, kept separately}
\label{sec:layers}
This split into a stable, saved \emph{scale} and a live, unsaved \emph{zero} is the
one idea to hold on to. In full, the system keeps \textbf{three} things apart, each
in its own place --- and only the first two are ever saved:
\begin{center}
\begin{tabular}{@{}lll@{}}
\toprule
\textbf{What} & \textbf{What it is} & \textbf{Saved?} \\
\midrule
Active calibration & axis switches $+$ per-channel \emph{scales} & yes (durable) \\
Calibration profiles & a named library of saved calibrations & yes (durable) \\
Zero (tare) & the live per-channel zero point & no (live only) \\
\bottomrule
\end{tabular}
\end{center}
The \textbf{active calibration} is the scale set currently in force --- what you
produce in the wizard. A \textbf{profile} is a named snapshot of an active
calibration you can save and reload (Section~\ref{sec:profiles}), handy when you
swap between walls or sensor sets. The \textbf{zero} is the live tare
(Section~\ref{sec:zeroing}): cached in your browser, but never sent to the backend
and never part of a calibration or a profile.

\subsection{Why known weights --- and why three of them?}
We create known forces by hanging known masses. A mass $m$ produces a downward
gravitational force
\begin{equation}
  F = m\,g, \qquad g = 9.80665~\mathrm{m/s^2},
\end{equation}
so $10~\mathrm{kg}$ gives about $98~\mathrm{N}$ and $20~\mathrm{kg}$ about
$196~\mathrm{N}$. By recording the raw sensor value at several known weights we can
work out the two unknowns in Equation~\eqref{eq:line}.

The standard procedure uses \textbf{three weights: $0$, $10$ and $20$~kg} (you can
change these values in the app). Why three?
\begin{itemize}[leftmargin=1.4em]
  \item The \textbf{$0$~kg} measurement gives the fit its \textbf{zero reference}
        --- the reading with no load --- so the slope is measured from the right
        baseline. (This reference is used only to work out the scale; it is
        \emph{not} saved as the sensor's zero. Setting the live zero is a separate,
        on-demand step --- Section~\ref{sec:zeroing}.)
  \item A \textbf{single loaded weight} would already be enough to estimate the
        \textbf{scale} (slope) --- geometrically, two points define a straight
        line.
  \item The \textbf{third weight} is what makes the result trustworthy. With three
        points we can \emph{check} that the sensor really behaves like a straight
        line, and we can \emph{average out} measurement noise instead of trusting
        a single reading. More points $\Rightarrow$ a more reliable slope.
\end{itemize}
In short: one point anchors the baseline, the others fix and verify the slope ---
and the slope is the only thing calibration keeps.

\subsection{The method: a best-fit straight line}
We do not just connect dots --- we fit the \textbf{best straight line} through the
measured points using the method of \emph{least squares} (the line that sits as
close as possible to all points at once). Because we already know the zero point
from the $0$~kg reading, the line is forced through the origin once the offset is
subtracted. The slope of that line is the \textbf{scale}:
\begin{equation}
  \text{scale} \;=\; \frac{\sum_k x_k\,F_k}{\sum_k x_k^{2}},
  \qquad x_k = (\text{raw}_k - \text{offset}),
  \label{eq:slope}
\end{equation}
where $F_k$ is the known force at step $k$ and $x_k$ is the offset-corrected raw
reading. This is exactly the formula the app uses internally; you never compute it
by hand --- the wizard does it the moment you finish capturing the weights.

This slope is not arbitrary: it is the one that makes the line fit the points as
well as possible. We choose \textbf{scale} to minimise the total squared error
between the known forces and the line's predictions,
\begin{equation}
  S(\text{scale}) \;=\; \sum_k \bigl(F_k - \text{scale}\cdot x_k\bigr)^2 .
  \label{eq:lsq}
\end{equation}
Setting the derivative to zero, $\frac{\mathrm{d}S}{\mathrm{d}(\text{scale})}
= -2\sum_k x_k\bigl(F_k - \text{scale}\,x_k\bigr) = 0$, and solving for
\textbf{scale} gives Equation~\eqref{eq:slope}. Since $S$ is an upward-opening
parabola in \textbf{scale} (its second derivative $2\sum_k x_k^{2}$ is positive),
this critical point is the unique minimum --- the best-fit slope.

\begin{infobox}
\textbf{In one sentence:} calibration applies known weights to each sensor, fits a
straight line through the raw readings, and stores the line's \textbf{slope}
(\emph{scale}) so that every future reading can be turned into newtons. The line's
zero point is \emph{not} saved --- you set it live with \textsf{Zero Sensors}
(Section~\ref{sec:zeroing}).
\end{infobox}

% =====================================================================
\section{The wall is tilted --- and why that helps}
% =====================================================================

The climbing wall does not stand vertically; it \textbf{leans back} by an angle we
call $\theta$ (theta). This tilt is not a nuisance --- it is the key to a clever
shortcut.

Each board measures force along two perpendicular directions that matter here:
\begin{itemize}[leftmargin=1.4em]
  \item $X$ --- \emph{along the wall surface} (the ``in-plane'' or up-slope
        direction),
  \item $Z$ --- \emph{perpendicular to the wall surface} (the ``normal'' or
        out-of-wall direction),
  \item $Y$ --- sideways (left/right), handled separately.
\end{itemize}

When you hang a weight straight down on a wall that is tilted by $\theta$, gravity
no longer points purely along one axis. Instead the single vertical force $W=mg$
splits into two parts (Figure~\ref{fig:tilt}):
\begin{equation}
  \underbrace{W\cos\theta}_{\text{along the wall }(X)} \qquad\text{and}\qquad
  \underbrace{W\sin\theta}_{\text{out of the wall }(Z)} .
\end{equation}

\begin{figure}[H]
\centering
\begin{tikzpicture}[scale=1.0,>=latex]
  \def\th{24}
  % define the geometry first
  \coordinate (O) at (0,0);
  \coordinate (A) at (0,4);
  \coordinate (B) at ({4*sin(\th)},{4*cos(\th)});
  % a point on the wall where the weight hangs
  \coordinate (P) at ({2.2*sin(\th)},{2.2*cos(\th)});
  % vertical reference
  \draw[gray,dashed] (O) -- (A) node[above]{\small vertical};
  % wall, tilted by theta from vertical
  \draw[very thick,brown!70!black] (O) -- (B)
        node[above right]{\small wall surface};
  % angle theta between vertical and wall (small wedge, not the reflex angle)
  \draw pic["$\theta$",draw=black,angle radius=0.9cm,angle eccentricity=1.45]
        {angle=B--O--A};
  \fill (P) circle (1.6pt);
  % weight pulling straight down
  \draw[->,very thick,blue!70!black] (P) -- ($(P)+(0,-1.6)$)
        node[below]{\small $W=mg$};
  % along-wall component (down-slope) at P
  \draw[->,thick,red!70!black] (P) -- ($(P)+({-1.25*sin(\th)},{-1.25*cos(\th)})$)
        node[below left]{\small $W\cos\theta\ (X)$};
  % normal component (into/out of wall) at P
  \draw[->,thick,green!45!black] (P) -- ($(P)+({-1.0*cos(\th)},{1.0*sin(\th)})$)
        node[left]{\small $W\sin\theta\ (Z)$};
\end{tikzpicture}
\caption{On a wall tilted by $\theta$ from vertical, one hanging weight $W$ loads
both the along-wall axis ($X$, by $W\cos\theta$) and the out-of-wall axis ($Z$, by
$W\sin\theta$) at the same time.}
\label{fig:tilt}
\end{figure}

This is the so-called \textbf{single-hang method}: because the tilt automatically
spreads the load across both the $X$ and $Z$ axes, \textbf{one set of hangs
calibrates two axes at once}. You only have to tell the app the wall's tilt angle
$\theta$ so it knows how the weight was divided. (The third axis, $Y$ /
sideways, is calibrated with a separate left/right pull.)

\begin{warnbox}
\textbf{The angle $\theta$ must match the real wall.} $\theta$ is used only to
split a hang into its $X$ and $Z$ parts. If you type the wrong angle, the split is
wrong and the calibration comes out wrong. Measure the wall's lean-back before you
start (a phone inclinometer app is fine). At $\theta\approx 0^\circ$ a hang puts no
load on $Z$ at all, so the wall \emph{must} be tilted for this method to work.
\end{warnbox}

\subsection{The coordinate convention}
Throughout the system the directions are fixed as:
\begin{center}
\begin{tabular}{@{}cl@{}}
\toprule
\textbf{Axis} & \textbf{Positive direction} \\
\midrule
$+X$ & up (toward the top of the wall) \\
$+Y$ & to the right \\
$+Z$ & out of the wall (toward the climber) \\
\bottomrule
\end{tabular}
\end{center}
A climber simply hanging from a hold therefore registers a \emph{negative} $X$
(they pull down-slope), and a pull away from the wall registers a \emph{positive}
$Z$.

% =====================================================================
\section{Before you begin}
% =====================================================================

You will need:
\begin{itemize}[leftmargin=1.4em]
  \item the backend running and the dashboard open in a browser (see the
        \emph{User Guide} for how to start them);
  \item the four sensor boards \textbf{plugged in} (via the USB hub) and the
        dashboard \textbf{connected} (click \textsf{Connect Device} in the
        toolbar) --- calibration reads live data, so without a connection every
        capture will time out;
  \item \textbf{known weights}, e.g.\ $10~\mathrm{kg}$ and $20~\mathrm{kg}$ plates,
        plus the ability to hang/pull them cleanly on each board;
  \item the wall's \textbf{tilt angle $\theta$}, measured in degrees from vertical.
\end{itemize}

\begin{infobox}
You calibrate \textbf{four boards}, and for each board you calibrate the
\textbf{hang ($X$ \& $Z$)} and the \textbf{$Y$ (sideways)} direction. That is the
full procedure, and it sets the \textbf{scales} only. Axes you skip keep rough
factory-default values and may report inaccurate forces. Zeroing (the tare) is a
separate one-click step done from the main toolbar and needs no weights ---
Section~\ref{sec:zeroing}.
\end{infobox}

\newpage

% =====================================================================
\section{Walkthrough: calibrating step by step}
% =====================================================================

Open the wizard by clicking \textsf{Calibrate} in the toolbar's
\emph{sensor-calibration} tray (alongside \textsf{Zero Sensors} and
\textsf{Profiles}). The wizard sets the \textbf{scales}, and walks you through five
stages:
\[
\text{Overview}\;\rightarrow\;\text{Set angle }\theta\;\rightarrow\;
\text{Pick a board}\;\rightarrow\;\text{Pick what to calibrate}\;\rightarrow\;
\text{Capture \& save}.
\]

\subsection{Stage 1 --- Overview}
The first screen (Figure~\ref{fig:calib_confirm}) summarises the procedure and
reminds you that hanging one weight calibrates both $X$ and $Z$. If you see an
amber ``No device connected'' banner, plug in and connect the device first. Click
\textsf{Start Calibration} to continue.

\screenshot{calib_confirm}{0.82}{Stage 1 --- the overview screen explains the
procedure before you begin.}

\subsection{Stage 2 --- Set the wall angle $\theta$}
Enter the wall's true lean-back from vertical (Figure~\ref{fig:calib_angle}). This
is the \emph{only} place the calibration angle can be set, on purpose: it keeps the
saved $X/Z$ calibration and the angle from ever drifting apart. The default is
$16^\circ$. Type your measured value and click \textsf{Next --- pick a board}.
The dashboard's \textsf{View in world frame} dialog asks for a tilt as well, but
that is a \emph{separate, display-only} angle --- pre-filled from this one and
never written back to it.

\screenshot{calib_angle}{0.82}{Stage 2 --- enter the measured wall tilt $\theta$.
It tells the app how a hang divides between the $X$ and $Z$ axes.}

\subsection{Stage 3 --- Pick a board}
A 2$\times$2 grid represents the wall as you face it: hands on top, feet on the
bottom (Figure~\ref{fig:calib_boards}). Each board shows three little badges, $X$,
$Y$ and $Z$. A \textbf{green check} means that axis is already calibrated; a plain
letter means it still uses factory defaults. Click the board you want to work on.

\screenshot{calib_boards}{0.82}{Stage 3 --- choose a board. Green checks mark axes
that are already calibrated.}

\subsection{Stage 4 --- Pick what to calibrate}
Now choose \emph{what} to calibrate on this board
(Figure~\ref{fig:calib_mode}):
\begin{itemize}[leftmargin=1.4em]
  \item \textbf{Hang $\rightarrow X$ \& $Z$} --- the single-hang method. You hang
        weights straight down and the app calibrates both the along-wall ($X$) and
        out-of-wall ($Z$) axes together. (Disabled if $\theta\approx 0$.)
  \item \textbf{$Y$ --- sideways} --- you pull the weight left or right to
        calibrate the sideways axis. The app knows that left-side boards are pulled
        left and flips the sign automatically.
\end{itemize}

\screenshot{calib_mode}{0.82}{Stage 4 --- calibrate the hang ($X$ \& $Z$) or the
sideways ($Y$) direction for the selected board.}

\subsection{Stage 5 --- Capture the weights and save}
This is where the measurements happen (Figure~\ref{fig:calib_steps}). You will see
three rows for $0$, $10$ and $20$~kg (these numbers are editable). For each row:
\begin{enumerate}[leftmargin=1.6em]
  \item Apply the stated weight. For a \textbf{hang}, let the weight hang
        \emph{straight down} and stay still. For \textbf{$Y$}, pull it cleanly to
        the side.
  \item Wait for it to stop moving, then click \textsf{Capture}. The app averages
        about 1.5~seconds of live data ($\sim$150 readings) to get a clean value.
  \item Repeat for all three weights. The $0$~kg step (nothing applied) gives the
        fit its zero reference --- it is used to work out the scale, not saved as
        the sensor's zero.
\end{enumerate}

A live \textbf{``Steady / Swinging''} indicator tells you whether the weight is
still moving. If you capture while it swings, the row is flagged ``swung'' --- you
are never blocked, but a steadier capture is more accurate, so re-capture once it
settles.

Once all three weights are captured, the app instantly computes and displays the
resulting \textbf{scale} (for a hang, one for $X$ and one for $Z$) --- no offset is
shown or stored. Review the green result panel, then click \textsf{Save}. The
calibration takes effect immediately and is stored permanently (see
Section~\ref{sec:storage}). You are returned to the board's mode picker so you can
calibrate its other axis, or go back and pick another board.

\screenshot{calib_steps}{0.86}{Stage 5 --- apply each known weight and click
\textsf{Capture}. The app fits the line and shows the computed scale before you
save.}

\subsection{Finishing}
Click \textsf{Done} when every board and axis is calibrated. If any axis is still
on factory defaults, the app warns you and lists exactly which ones, so nothing is
forgotten. There is also an \textsf{Exit \& discard} option that rolls back every
change you made during the session (including $\theta$), and a \textsf{Reset to
defaults} button at the bottom that clears \emph{all} calibration.

\begin{infobox}
Two things worth knowing about the wizard. The \textbf{green ``calibrated'' checks}
start empty every time you open it and are filled only by saves you make \emph{in
that session} --- they track this sitting, not your whole history. And the footer's
\textsf{Save / Load Profile} button opens the profile manager
(Section~\ref{sec:profiles}) without leaving the wizard.
\end{infobox}

% =====================================================================
\section{Zeroing the sensors (tare)}
\label{sec:zeroing}
% =====================================================================

Calibration fixes the \emph{scale}; \textbf{zeroing} fixes the \emph{zero point}.
Because the zero drifts, you re-set it whenever you start measuring --- it takes one
click, needs no weights, and never touches your calibration.

\screenshot{zero_sensors}{0.8}{The sensor-calibration tray: \textsf{Zero Sensors}
with its live status chip (\textsf{Zeroed $X$s ago} and a freshness dot), next to
\textsf{Calibrate} and \textsf{Profiles}.}

\paragraph{How to zero.}
\begin{enumerate}[leftmargin=1.6em]
  \item Make sure the holds are \textbf{unloaded} (nothing hanging, nobody on the
        wall) and the device is connected and streaming.
  \item Click \textsf{Zero Sensors} in the toolbar's sensor-calibration tray. The
        app averages about $1.5$~seconds of the most recent readings ($\sim$150
        samples) and installs that as the new zero \emph{instantly} --- a recording
        in progress never pauses.
\end{enumerate}

\paragraph{Reading the status.} Next to the button a small chip reports the state:
\textsf{Not zeroed yet} before you have zeroed, then \textsf{Zeroed $X$s ago} with a
coloured dot --- \textcolor{green!55!black}{green} under a minute,
\textcolor{warn}{amber} under five, red beyond --- so you can tell at a glance
whether the zero is fresh. The small $\circlearrowleft$ beside it clears the zero
and returns to raw, un-tared readings.

\begin{infobox}
The zero is \textbf{cached in your browser}, so it survives a page reload and you do
not have to re-zero after refreshing. But it is \textbf{never sent to the backend}
and is \textbf{never part of a calibration or profile}: it is the one value the
system deliberately keeps live rather than saved.
\end{infobox}

\begin{warnbox}
\textbf{Zeroing needs a live device.} The button stays disabled until about
$20$~samples have streamed in, and \textbf{demo mode does not feed it} (demo data
bypasses the sensor pipeline, the same limitation as the wizard's live readout).
Connect a real device to zero.
\end{warnbox}

\paragraph{When to zero.} At the start of each session, and any time the unloaded
reading looks off or the staleness dot turns red. You do \emph{not} need to
recalibrate in order to re-zero --- keeping the two apart is exactly the point.

% =====================================================================
\section{Calibration profiles}
\label{sec:profiles}
% =====================================================================

A \textbf{profile} is a named, saved copy of a full calibration (all the axis
switches and scales --- never the zero). Profiles let you keep more than one
calibration on hand: one per wall, one per set of sensors, or a known-good backup
before you experiment. They live on the backend and survive restarts.

Open the manager from the \textsf{Profiles} button in the toolbar, or from
\textsf{Save / Load Profile} in the calibration wizard's footer --- both open the
same \textsf{Calibration Profiles} panel.

\screenshot{calib_profiles}{0.72}{The \textsf{Calibration Profiles} panel: name and
\textsf{Save as new} at the top, then each saved profile with \textsf{Load},
duplicate, rename and delete.}

\paragraph{Saving.} Type a name under \textsf{Save current calibration} and click
\textsf{Save as new} to store whatever calibration is currently active. If a profile
is already active, an \textsf{Update ``\emph{name}''} button also appears to write
your latest changes back to it.

\paragraph{Managing.} Each saved profile offers five actions:
\begin{itemize}[leftmargin=1.4em]
  \item \textsf{Load} --- apply it as the active calibration (it takes effect and is
        saved immediately);
  \item \textbf{Recalibrate} (ruler icon) --- load the profile \emph{and} open the
        calibration wizard on it in one click. This is the only way to change a
        profile's values: they are always measured with real weights in the wizard,
        never typed in by hand;
  \item \textbf{Duplicate} --- copy it under a new name;
  \item \textbf{Rename} --- change the name only (the values stay untouched);
  \item \textbf{Delete} --- guarded by a two-step confirmation.
\end{itemize}
The active profile is marked \textsf{Active}, or \textsf{Active \textperiodcentered\
modified} once you have changed the live calibration without saving it back.

\begin{infobox}
\textbf{Editing a profile safely.} Click the profile's \textbf{Recalibrate} (ruler)
button --- it becomes the active calibration and the wizard opens on it. A banner at
the top of the wizard always tells you \emph{which} calibration you are changing
(``loaded from profile `\emph{name}'\,'') and warns once your changes differ from the
saved profile. Recalibrate the boards you want --- each \textsf{Save} updates the
live calibration --- then click \textsf{Update ``\emph{name}''} to write the changes
back. \textsf{Duplicate} first if you want to keep the original untouched.
\end{infobox}

% =====================================================================
\section{Where it is all stored}
\label{sec:storage}
% =====================================================================

The three things from Section~\ref{sec:layers} are kept --- or deliberately not
kept --- in different places:
\begin{itemize}[leftmargin=1.4em]
  \item \textbf{The active calibration} (axis switches $+$ scales). Pressing
        \textsf{Save} in the wizard writes it in three places at once: the
        \textbf{live program} (the next reading is already converted correctly), the
        browser's \textbf{local storage} (a reload keeps it even with no backend),
        and a file on the backend, \texttt{calibration\_settings.json} (the
        authoritative copy, reloaded automatically next time). Only scales and axis
        switches go here --- \textbf{no zero}.
  \item \textbf{Profiles} live in \texttt{calibration\_profiles.json} on the
        backend, so your named library survives restarts.
  \item \textbf{The zero (tare)} is cached \emph{only} in the browser's local
        storage. It is never persisted to the backend and never part of a
        calibration or profile.
\end{itemize}
Because the scales are a physical property of the sensors, you normally calibrate
\textbf{once} per board; you re-zero often, but that costs nothing and saves
nothing.

\begin{infobox}
\textbf{Do I need to recalibrate if the wall angle changes?} Generally no. The
scales are intrinsic to the sensors; $\theta$ only affects how a hang is split at
calibration time and how forces are \emph{displayed}. Re-running the
\textbf{Hang $\rightarrow X$ \& $Z$} step after a large angle change is only a
precaution against real-world cross-axis effects. The $Y$ axis and the overall
force magnitude are unaffected either way.
\end{infobox}

% =====================================================================
\section{Troubleshooting}
% =====================================================================

\begin{center}
\begin{tabular}{@{}p{0.42\linewidth}p{0.52\linewidth}@{}}
\toprule
\textbf{Symptom} & \textbf{What to do} \\
\midrule
``No sensor data received'' when capturing &
Make sure the backend is running, the sensor boards are plugged in, and the
dashboard is connected (\textsf{Connect Device}). Captures need live data and time
out after 8~seconds without it. \\[4pt]
The \textsf{Zero Sensors} button is greyed out &
No live data is streaming. Plug in the boards, click \textsf{Connect Device} and
wait a moment --- zeroing needs about $20$~samples, and demo mode cannot zero. \\[4pt]
Unloaded readings sit at a non-zero value &
Re-zero: unload the holds and click \textsf{Zero Sensors}. A red staleness dot means
the old zero has simply drifted. \\[4pt]
The \textbf{Hang} option is greyed out &
The wall angle $\theta$ is set to roughly $0^\circ$. A vertical hang then puts no
load on $Z$. Go back to Stage~2 and enter the real (non-zero) tilt. \\[4pt]
A captured row says ``swung $\pm N$'' &
The weight was still moving. Let it settle and \textsf{Re-capture} for a cleaner
value. \\[4pt]
Forces look wrong after calibration &
First re-zero (unloaded) with \textsf{Zero Sensors}. If they are still off, check
that $\theta$ matched the real wall and that the weights you entered match the
plates you used. \\
\bottomrule
\end{tabular}
\end{center}

\vfill
\begin{center}\small\itshape\color{muted}
Screenshots were taken from the dashboard running in demo mode; on a connected
system the ``No device connected'' banner does not appear.
\end{center}

\end{document}
